\documentclass[10pt,a4paper]{amsart}
\usepackage{amsmath,amssymb,amsthm,graphicx}
\usepackage{verbatim}
\newenvironment{reference}{\comment}{\endcomment}
%\newenvironment{reference}{}{}
\newenvironment{slogan}{\comment}{\endcomment}
\newenvironment{history}{\comment}{\endcomment}

% For commutative diagrams we use Xy-pic
\usepackage[all]{xy}

% We use 2cell for 2-commutative diagrams.
\xyoption{2cell}
% Theorem environments.
%
\theoremstyle{plain}
\newtheorem{theorem}[subsection]{Theorem}
\newtheorem{proposition}[subsection]{Proposition}
\newtheorem{lemma}[subsection]{Lemma}

\theoremstyle{definition}
\newtheorem{definition}[subsection]{Definition}
\newtheorem{example}[subsection]{Example}
\newtheorem{exercise}[subsection]{Exercise}
\newtheorem{situation}[subsection]{Situation}

\theoremstyle{remark}
\newtheorem{remark}[subsection]{Remark}
\newtheorem{remarks}[subsection]{Remarks}

\numberwithin{equation}{subsection}

% Macros
%
\def\lim{\mathop{\mathrm{lim}}\nolimits}
\def\colim{\mathop{\mathrm{colim}}\nolimits}
\def\Spec{\mathop{\mathrm{Spec}}}
\def\Hom{\mathop{\mathrm{Hom}}\nolimits}
\def\Ext{\mathop{\mathrm{Ext}}\nolimits}
\def\SheafHom{\mathop{\mathcal{H}\!\mathit{om}}\nolimits}
\def\SheafExt{\mathop{\mathcal{E}\!\mathit{xt}}\nolimits}
\def\Sch{\mathit{Sch}}
\def\Mor{\mathop{\mathrm{Mor}}\nolimits}
\def\Ob{\mathop{\mathrm{Ob}}\nolimits}
\def\Sh{\mathop{\mathit{Sh}}\nolimits}
\def\NL{\mathop{N\!L}\nolimits}
\def\CH{\mathop{\mathrm{CH}}\nolimits}
\def\proetale{{pro\text{-}\acute{e}tale}}
\def\etale{{\acute{e}tale}}
\def\QCoh{\mathit{QCoh}}
\def\Ker{\mathop{\mathrm{Ker}}}
\def\Im{\mathop{\mathrm{Im}}}
\def\Coker{\mathop{\mathrm{Coker}}}
\def\Coim{\mathop{\mathrm{Coim}}}

% Boxtimes
%
\DeclareMathSymbol{\boxtimes}{\mathbin}{AMSa}{"02}

%
% Macros for moduli stacks/spaces
%
\def\QCohstack{\mathcal{QC}\!\mathit{oh}}
\def\Cohstack{\mathcal{C}\!\mathit{oh}}
\def\Spacesstack{\mathcal{S}\!\mathit{paces}}
\def\Quotfunctor{\mathrm{Quot}}
\def\Hilbfunctor{\mathrm{Hilb}}
\def\Curvesstack{\mathcal{C}\!\mathit{urves}}
\def\Polarizedstack{\mathcal{P}\!\mathit{olarized}}
\def\Complexesstack{\mathcal{C}\!\mathit{omplexes}}
% \Pic is the operator that assigns to X its picard group, usage \Pic(X)
% \Picardstack_{X/B} denotes the Picard stack of X over B
% \Picardfunctor_{X/B} denotes the Picard functor of X over B
\def\Pic{\mathop{\mathrm{Pic}}\nolimits}
\def\Picardstack{\mathcal{P}\!\mathit{ic}}
\def\Picardfunctor{\mathrm{Pic}}
\def\Deformationcategory{\mathcal{D}\!\mathit{ef}}
\usepackage{cleveref}
\crefname{originallemma}{original lemma}{original lemmas}
\Crefname{originallemma}{Original lemma}{Original lemmas}
\makeatletter
\newlabel{categories-lemma-yoneda}{{001P}{source}{}{sourcepointer-1}{}}
\newlabel{categories-lemma-yoneda@cref}{{[originallemma][0][]001P}{[1][0][]source}{}{}{}}
\makeatother
\crefname{originallemma}{original lemma}{original lemmas}
\Crefname{originallemma}{Original lemma}{Original lemmas}
\makeatletter
\newlabel{groupoids-lemma-colimit-kappa}{{077U}{source}{}{sourcepointer-2}{}}
\newlabel{groupoids-lemma-colimit-kappa@cref}{{[originallemma][0][]077U}{[1][0][]source}{}{}{}}
\makeatother
\crefname{originallemma}{original lemma}{original lemmas}
\Crefname{originallemma}{Original lemma}{Original lemmas}
\makeatletter
\newlabel{groupoids-lemma-set-of-iso-classes}{{077T}{source}{}{sourcepointer-3}{}}
\newlabel{groupoids-lemma-set-of-iso-classes@cref}{{[originallemma][0][]077T}{[1][0][]source}{}{}{}}
\makeatother
\crefname{originallemma}{original lemma}{original lemmas}
\Crefname{originallemma}{Original lemma}{Original lemmas}
\makeatletter
\newlabel{injectives-lemma-grothendieck-brown}{{07D7}{source}{}{sourcepointer-4}{}}
\newlabel{injectives-lemma-grothendieck-brown@cref}{{[originallemma][0][]07D7}{[1][0][]source}{}{}{}}
\makeatother
\crefname{originallemma}{original lemma}{original lemmas}
\Crefname{originallemma}{Original lemma}{Original lemmas}
\makeatletter
\newlabel{injectives-lemma-grothendieck-products}{{07D8}{source}{}{sourcepointer-5}{}}
\newlabel{injectives-lemma-grothendieck-products@cref}{{[originallemma][0][]07D8}{[1][0][]source}{}{}{}}
\makeatother
\crefname{originalsection}{original section}{original sections}
\Crefname{originalsection}{Original section}{Original sections}
\makeatletter
\newlabel{injectives-section-grothendieck-conditions}{{079A}{source}{}{sourcepointer-6}{}}
\newlabel{injectives-section-grothendieck-conditions@cref}{{[originalsection][0][]079A}{[1][0][]source}{}{}{}}
\makeatother
\crefname{originaltheorem}{original theorem}{original theorems}
\Crefname{originaltheorem}{Original theorem}{Original theorems}
\makeatletter
\newlabel{injectives-theorem-injective-embedding-grothendieck}{{079H}{source}{}{sourcepointer-7}{}}
\newlabel{injectives-theorem-injective-embedding-grothendieck@cref}{{[originaltheorem][0][]079H}{[1][0][]source}{}{}{}}
\makeatother
\crefname{originalreference}{original reference}{original references}
\Crefname{originalreference}{Original reference}{Original references}
\makeatletter
\newlabel{lemma-properties-quasi-coherent}{{03M1}{source}{}{sourcepointer-8}{}}
\newlabel{lemma-properties-quasi-coherent@cref}{{[originalreference][0][]03M1}{[1][0][]source}{}{}{}}
\makeatother
\crefname{originalproposition}{original proposition}{original propositions}
\Crefname{originalproposition}{Original proposition}{Original propositions}
\makeatletter
\newlabel{properties-proposition-coherator}{{077P}{source}{}{sourcepointer-9}{}}
\newlabel{properties-proposition-coherator@cref}{{[originalproposition][0][]077P}{[1][0][]source}{}{}{}}
\makeatother
\crefname{originalreference}{original reference}{original references}
\Crefname{originalreference}{Original reference}{Original references}
\makeatletter
\newlabel{proposition-quasi-coherent}{{03M3}{source}{}{sourcepointer-10}{}}
\newlabel{proposition-quasi-coherent@cref}{{[originalreference][0][]03M3}{[1][0][]source}{}{}{}}
\makeatother
\crefname{originallemma}{original lemma}{original lemmas}
\Crefname{originallemma}{Original lemma}{Original lemmas}
\makeatletter
\newlabel{sites-modules-lemma-limits-colimits}{{03DB}{source}{}{sourcepointer-11}{}}
\newlabel{sites-modules-lemma-limits-colimits@cref}{{[originallemma][0][]03DB}{[1][0][]source}{}{}{}}
\makeatother
\title{Quasi-coherent sheaves on an algebraic space form a Grothendieck abelian category and the inclusion has a right adjoint}
\author{The Stacks Project Authors}
\date{stacks-project:spaces-properties.tex:local-snapshot-20261012; GNU FDL 1.2 or later}
\begin{document}
\maketitle
\noindent\textit{Editorial note.} All delivered mathematical text is the author's own, copied from the pinned Stacks Project chapter file spaces-properties.tex; no statement or proof step was rewritten, reordered or removed.\par\medskip
\noindent\textit{Editorial note.} The counted proof is the author's complete argument: the verification of the Grothendieck conditions through the closure properties of quasi-coherent sheaves, the construction of a generator from the groupoid presentation using the boundedness of kappa-generated compatible systems, and the construction of the coherator by representing a contravariant Hom functor.\par\medskip
\noindent\textit{Editorial note.} The label proposition-coherator names the counted statement, and the proof environment immediately following it in the source is delivered verbatim as its proof.\par\medskip
\noindent\textit{Editorial note.} The counted statement is carried by the author's own proposition environment in the chapter source.\par\medskip
\noindent\textit{Editorial note.} Pre-open gate: python operations/agent-authoring/tools/precard\_gate.py --source spaces-properties.tex --label proposition-coherator --start 6023 --end 6105 --local-id 28045 returned RESULT=OK (L28045 free, occupancy FREE, public interval clear). Candidate taken from sources/stacks-spaces-properties/candidates-halfA.json.\par\medskip
\noindent\textit{Editorial note.} The proof opens with the sentence 'This proof is a repeat of the proof in the case of schemes ... We advise the reader to read that proof first.' — guidance about an analogous earlier proof, while the proof environment delivers the full argument for the present case, so the statement is counted and the sentence is recorded here verbatim. The statement also carries a footnote naming the adjoint, which is delivered verbatim.\par\medskip
\noindent\textit{Editorial note.} The statement and proof are copied from the pinned chapter file; the cited inputs are kept as external pointers with their canonical Stacks tags.\par\medskip
\noindent\textit{Allowed standard prerequisites (editorial).}\begin{itemize}
\item Closure properties of the category of quasi-coherent sheaves. Same chapter, the lemma that the category of quasi-coherent sheaves on an algebraic space is stable under colimits, kernels, cokernels, tensor products and internal homs. Supplies the existence of colimits and their agreement with colimits in all modules.
\item Equivalence between quasi-coherent sheaves on the space and compatible systems on a presentation. Same chapter, the proposition identifying quasi-coherent sheaves on an algebraic space with the compatible systems of quasi-coherent sheaves on an etale groupoid scheme presenting it. Supplies the construction of the quasi-coherent sheaves on the space corresponding to the chosen kappa-generated systems.
\item Boundedness and generators for compatible systems over a groupoid scheme. Groupoids, the lemma producing a cardinal bounding the generated submodules and the lemma bounding the set of isomorphism classes of generators. Supply the cardinal kappa and the collection of kappa-generated compatible systems used to build a generator of the category.
\item Brown representability and the Grothendieck abelian formalism. Injectives, the section on Grothendieck conditions together with the Brown representability lemma and the existence theorem for injective hulls. Supply the representability of the contravariant functor defining the coherator, and the consequences about limits and enough injectives.
\end{itemize}\medskip
\ifdefined\hypertarget\hypertarget{sourcepointer-1}{}\fi
\noindent\textit{Source reference.} Original lemma 001P: sources/stacks-chow/tags.txt, source lines 70--70. Prerequisite pointer invoked at proof line(s) 6096; it names a standard fact of a sibling chapter that the counted argument uses, and every step of the counted proof built on it is retained verbatim in the delivered body.\par\medskip
\ifdefined\hypertarget\hypertarget{sourcepointer-2}{}\fi
\noindent\textit{Source reference.} Original lemma 077U: sources/stacks-chow/tags.txt, source lines 9214--9214. Prerequisite pointer invoked at proof line(s) 6062; it names a standard fact of a sibling chapter that the counted argument uses, and every step of the counted proof built on it is retained verbatim in the delivered body.\par\medskip
\ifdefined\hypertarget\hypertarget{sourcepointer-3}{}\fi
\noindent\textit{Source reference.} Original lemma 077T: sources/stacks-chow/tags.txt, source lines 9213--9213. Prerequisite pointer invoked at proof line(s) 6066; it names a standard fact of a sibling chapter that the counted argument uses, and every step of the counted proof built on it is retained verbatim in the delivered body.\par\medskip
\ifdefined\hypertarget\hypertarget{sourcepointer-4}{}\fi
\noindent\textit{Source reference.} Original lemma 07D7: sources/stacks-chow/tags.txt, source lines 9402--9402. Prerequisite pointer invoked at proof line(s) 6091; it names a standard fact of a sibling chapter that the counted argument uses, and every step of the counted proof built on it is retained verbatim in the delivered body.\par\medskip
\ifdefined\hypertarget\hypertarget{sourcepointer-5}{}\fi
\noindent\textit{Source reference.} Original lemma 07D8: sources/stacks-chow/tags.txt, source lines 9403--9403. Prerequisite pointer invoked at proof line(s) 6079; it names a standard fact of a sibling chapter that the counted argument uses, and every step of the counted proof built on it is retained verbatim in the delivered body.\par\medskip
\ifdefined\hypertarget\hypertarget{sourcepointer-6}{}\fi
\noindent\textit{Source reference.} Original section 079A: sources/stacks-chow/tags.txt, source lines 9265--9265. Prerequisite pointer invoked at proof line(s) 6050; it names a standard fact of a sibling chapter that the counted argument uses, and every step of the counted proof built on it is retained verbatim in the delivered body.\par\medskip
\ifdefined\hypertarget\hypertarget{sourcepointer-7}{}\fi
\noindent\textit{Source reference.} Original theorem 079H: sources/stacks-chow/tags.txt, source lines 9272--9272. Prerequisite pointer invoked at proof line(s) 6078; it names a standard fact of a sibling chapter that the counted argument uses, and every step of the counted proof built on it is retained verbatim in the delivered body.\par\medskip
\ifdefined\hypertarget\hypertarget{sourcepointer-8}{}\fi
\noindent\textit{Source reference.} Original reference 03M1: sources/stacks-spaces-properties/pinned-original-source.tex, source lines 5658--5658. Prerequisite pointer invoked at proof line(s) 6051; it names a standard fact of the same chapter that the counted argument uses, and every step of the counted proof built on it is retained verbatim in the delivered body.\par\medskip
\ifdefined\hypertarget\hypertarget{sourcepointer-9}{}\fi
\noindent\textit{Source reference.} Original proposition 077P: sources/stacks-chow/tags.txt, source lines 9209--9209. Prerequisite pointer invoked at proof line(s) 6044; it names a standard fact of a sibling chapter that the counted argument uses, and every step of the counted proof built on it is retained verbatim in the delivered body.\par\medskip
\ifdefined\hypertarget\hypertarget{sourcepointer-10}{}\fi
\noindent\textit{Source reference.} Original reference 03M3: sources/stacks-spaces-properties/pinned-original-source.tex, source lines 5944--5944. Prerequisite pointer invoked at proof line(s) 6070; it names a standard fact of the same chapter that the counted argument uses, and every step of the counted proof built on it is retained verbatim in the delivered body.\par\medskip
\ifdefined\hypertarget\hypertarget{sourcepointer-11}{}\fi
\noindent\textit{Source reference.} Original lemma 03DB: sources/stacks-chow/tags.txt, source lines 4414--4414. Prerequisite pointer invoked at proof line(s) 6054; it names a standard fact of a sibling chapter that the counted argument uses, and every step of the counted proof built on it is retained verbatim in the delivered body.\par\medskip
\begin{proposition}
\label{proposition-coherator}
Let $S$ be a scheme. Let $X$ be an algebraic space over $S$.
\begin{enumerate}
\item The category $\QCoh(\mathcal{O}_X)$ is a Grothendieck
abelian category. Consequently, $\QCoh(\mathcal{O}_X)$
has enough injectives and all limits.
\item The inclusion functor
$\QCoh(\mathcal{O}_X) \to \textit{Mod}(\mathcal{O}_X)$
has a right adjoint\footnote{This functor is sometimes called
the {\it coherator}.}
$$
Q : \textit{Mod}(\mathcal{O}_X) \longrightarrow \QCoh(\mathcal{O}_X)
$$
such that for every quasi-coherent sheaf $\mathcal{F}$ the adjunction mapping
$Q(\mathcal{F}) \to \mathcal{F}$ is an isomorphism.
\end{enumerate}
\end{proposition}

\begin{proof}
This proof is a repeat of the proof in the case of schemes, see
Properties, Proposition \ref{properties-proposition-coherator}.
We advise the reader to read that proof first.

\medskip\noindent
Part (1) means $\QCoh(\mathcal{O}_X)$ (a) has all colimits,
(b) filtered colimits are exact, and (c) has a generator, see
Injectives, Section \ref{injectives-section-grothendieck-conditions}.
By Lemma \ref{lemma-properties-quasi-coherent}
colimits in $\QCoh(\mathcal{O}_X)$ exist and agree
with colimits in $\textit{Mod}(\mathcal{O}_X)$. By
Modules on Sites, Lemma \ref{sites-modules-lemma-limits-colimits}
filtered colimits are exact. Hence (a) and (b) hold.

\medskip\noindent
To construct a generator, choose a presentation $X = U/R$ so that
$(U, R, s, t, c)$ is an
\'etale groupoid scheme and in particular $s$ and $t$ are flat morphisms
of schemes. Pick a cardinal $\kappa$ as in
Groupoids, Lemma \ref{groupoids-lemma-colimit-kappa}.
Pick a collection $(\mathcal{E}_t, \alpha_t)_{t \in T}$ of
$\kappa$-generated quasi-coherent modules on
$(U, R, s, t, c)$ as in
Groupoids, Lemma \ref{groupoids-lemma-set-of-iso-classes}.
Let $\mathcal{F}_t$ be the quasi-coherent module on $X$ which
corresponds to the quasi-coherent module $(\mathcal{E}_t, \alpha_t)$ via
the equivalence of categories of
Proposition \ref{proposition-quasi-coherent}.
Then we see that every quasi-coherent module $\mathcal{H}$ is the
directed colimit of its quasi-coherent submodules which are isomorphic
to one of the $\mathcal{F}_t$. Thus $\bigoplus_t \mathcal{F}_t$ is
a generator of $\QCoh(\mathcal{O}_X)$ and we conclude that (c) holds.
The assertions on limits and injectives hold in any
Grothendieck abelian category, see
Injectives, Theorem
\ref{injectives-theorem-injective-embedding-grothendieck} and
Lemma \ref{injectives-lemma-grothendieck-products}.

\medskip\noindent
Proof of (2). To construct $Q$ we use the following general procedure.
Given an object $\mathcal{F}$ of $\textit{Mod}(\mathcal{O}_X)$
we consider the functor
$$
\QCoh(\mathcal{O}_X)^{opp} \longrightarrow \textit{Sets},\quad
\mathcal{G} \longmapsto \Hom_X(\mathcal{G}, \mathcal{F})
$$
This functor transforms colimits into limits,
hence is representable, see
Injectives, Lemma \ref{injectives-lemma-grothendieck-brown}.
Thus there exists a quasi-coherent sheaf $Q(\mathcal{F})$
and a functorial isomorphism
$\Hom_X(\mathcal{G}, \mathcal{F}) = \Hom_X(\mathcal{G}, Q(\mathcal{F}))$
for $\mathcal{G}$ in $\QCoh(\mathcal{O}_X)$. By the Yoneda lemma
(Categories, Lemma \ref{categories-lemma-yoneda})
the construction $\mathcal{F} \leadsto Q(\mathcal{F})$ is
functorial in $\mathcal{F}$. By construction $Q$ is a right
adjoint to the inclusion functor.
The fact that $Q(\mathcal{F}) \to \mathcal{F}$ is an isomorphism
when $\mathcal{F}$ is quasi-coherent is a formal consequence of the fact
that the inclusion functor
$\QCoh(\mathcal{O}_X) \to \textit{Mod}(\mathcal{O}_X)$
is fully faithful.
\end{proof}

\end{document}
